\documentclass[11pt,english]{article}
\bibliographystyle{jf}
\usepackage[longnamesfirst]{natbib}
%\usepackage[T1]{fontenc}
\usepackage{fontenc}
\usepackage[utf8]{inputenc}
\usepackage[bottom]{footmisc}
\usepackage{tabularx}
\usepackage{booktabs}
\usepackage{amsmath,enumerate}
\usepackage{amsthm}
\usepackage{dsfont}
\usepackage{amssymb}
\usepackage{mathtools}
\usepackage{mathrsfs}
\usepackage{graphicx}
\usepackage{natbib}
\usepackage[labelfont=bf,justification=justified,font=small,skip=2pt,labelsep=period]{caption}
%\usepackage[labelfont=bf,justification=justified,labelsep=period]{caption}
\usepackage{subcaption}
\usepackage{setspace}
\usepackage{arydshln}
\usepackage{chngpage}
\usepackage{float,appendix}
\usepackage{afterpage}
\usepackage{verbatim}
\usepackage{indentfirst}
\usepackage[margin=1in]{geometry}
\usepackage[usenames, dvipsnames]{color}
\usepackage{hyperref}
\usepackage{adjustbox}
% \usepackage{longtable}
\usepackage{booktabs, makecell, longtable,ltxtable}
\usepackage{lipsum}
\usepackage{pdflscape}
\usepackage{siunitx}
\usepackage{chngcntr}
\usepackage{lipsum}
\usepackage{tikz}
\usepackage{pdflscape}
\usepackage[normalem]{ulem}
\usepackage{chngcntr,apptools}
\AtAppendix{\counterwithin{lemma}{section}}
\usepackage{siunitx}
\usepackage{tabularx}
\sisetup{table-format=-1.3, table-space-text-post={**}}
\usepackage{xr}
\usepackage{etoolbox}
\usetikzlibrary{positioning}
\newcommand\fnsep{\textsuperscript{,}}
\usepackage{tikz}
\usepackage{pdflscape}
\usepackage[normalem]{ulem}
\usepackage{chngcntr,apptools}
\AtAppendix{\counterwithin{lemma}{section}}
\usepackage{siunitx}
\usepackage{tabularx}
\usepackage{xr}
\usepackage{etoolbox}
\usetikzlibrary{positioning}

%% UCLA colors: 
% https://brand.ucla.edu/wp-content/uploads/ucla-color-palette-v10.pdf
%\definecolor{UCLA}{RGB}{50, 132, 191}   % UCLA Blue primary
%\definecolor{UCLA}{RGB}{52, 123, 173}   % UCLA Blue text only
\definecolor{UCLAblue}{RGB}{30, 75, 135} % secondary color	
\definecolor{UCLAgold}{RGB}{255, 232, 0} % UCLA Gold

\definecolor{usccardinal}{rgb}{0.6, 0.0, 0.0}
\definecolor{Matlabblue}{rgb}{0,0.447,0.741}
\definecolor{bleudefrance}{rgb}{0.19,0.55, 0.91}
\definecolor{cobalt}{rgb}{0.0, 0.28,0.67}
\definecolor{darkblue}{rgb}{0.0, 0.0,0.55}
\definecolor{darkcerulean}{rgb}{0.03,0.27, 0.49}
\definecolor{darkpowderblue}{rgb}{0.0,0.2, 0.6}
\definecolor{bleudefrance}{rgb}{0.19,0.55, 0.91}
%\usepackage[colorlinks=true,urlcolor=blue,linkcolor=blue,citecolor=blue]{hyperref}

\hypersetup{
	colorlinks,
	linkcolor={UCLAblue},
	citecolor={UCLAblue},
	urlcolor={UCLAblue},
}

\bibpunct{(}{)}{;}{a}{,}{,}
\newcommand{\E}{\mathbb{E}}
\newcommand{\var}{\mathrm{Var}}
\newcommand{\cov}{\mathrm{Cov}}
\setcounter{MaxMatrixCols}{20}
\newtheorem{theorem}{Theorem}
\newtheorem{lemma}{Lemma}
\newtheorem{defn}{Definition}
\newtheorem{prop}{Proposition}
\newtheorem{cor}{Corollary}
\urlstyle{same}
\usepackage{pgfplots} 
\pgfplotsset{compat=newest} 
\pgfplotsset{plot coordinates/math parser=false} 
% puts figures and tables at the end of the document
%\usepackage[nomarkers,nolists]{endfloat}               
% puts ONLY figures at the end of the document
%\usepackage[nomarkers,nolists,figuresonly]{endfloat}  
% allows more than one figure per page for endfloat
%\renewcommand{\efloatseparator}{\mbox{}}              
\usepackage{rotating}
% allows sideways tables and figures to be placed at the end of the document
%\DeclareDelayedFloatFlavor{sidewaystable}{table}
%\DeclareDelayedFloatFlavor{sidewaysfigure}{figure}

% \usepackage[nomarkers,nolists]{endfloat} 
% \renewcommand{\efloatseparator}{\mbox{}}              % allows more than one figure per page for endfloat

%\usepackage{mathpazo}
%\usepackage{mathpple}
%%%%%%%%%%%%%%%%%%%%%%%%%%%%%%%%
%\usepackage{palatino}
%%%%%%%%%%%%%%%%%%%%%%%%%%%%%%%%
\pgfplotsset{compat=newest} 
\pgfplotsset{plot coordinates/math parser=false} 
\newlength\figureheight 
\newlength\figurewidth 
\usepackage{mathtools}
\usetikzlibrary{calc}


%%%%%%%%%%%%%%%%%%%%%%%%%%
% Notes customization %
%%%%%%%%%%%%%%%%%%%%%%%%%%
\definecolor{webblue}{rgb}{0,0,0.6}
\definecolor{webgreen}{rgb}{0,.6,0}
\definecolor{webbrown}{rgb}{.6,0,0}


%- Commenting
\let\oldmarginpar\marginpar
\renewcommand\marginpar[1]{\-\oldmarginpar[\raggedleft\tiny #1]%
{\raggedright\tiny #1}}
\newcommand{\mknote}[1]{\marginpar{\textcolor{webblue}{\tiny{Mahyar: #1}}}}
\newcommand{\dsnote}[1]{\marginpar{\textcolor{webgreen}{\tiny{Dejanir: #1}}}}
\usepackage[en-US]{datetime2}
\usepackage{mathtools}
 \usepackage{xcolor}

\newcommand{\comments}[1]
{\par {\bfseries \color{blue} #1 \par}}

\usepackage[draft]{todonotes}

\title{\textbf{A Note on Portfolio Demand Elasticities}}
\author{ Dejanir Silva\thanks{Gies College of Business, University of Illinois at Urbana-Champaign, e-mail: \href{dejanir@illinois.edu}{dejanir@illinois.edu}.}}
\date{November 2020}
%\date{\today\\
%	\begin{small}
%		First draft: May 5, 2018
%\end{small}}
\begin{document}
	\maketitle
	
\section{A Lucas Economy with Labor Income}

Consider an economy where aggregate endowment follows a geometric Brownian motion
\begin{equation*}
    \frac{d Y_t}{Y_t} = \mu dt + \sigma d Z_t.
\end{equation*}

Investors have access to a riskless and a risky asset paying dividends $D_t = \psi Y_t$. Investors also receive non-business (labor) income $I_t = (1-\psi) Y_t$. In equilibrium, the interest rate $r$ and the risk premium $\pi$ will be constant. The price of the risky asset will be given by $P_t = \frac{D_t}{r+\pi -\mu}$. 

Investors have CRRA preferences and choose consumption $C_t$ and the share of financial wealth invested in the risky asset:
\begin{equation*}
    V(W, I) = \max_{[C_s, Q_s]} \mathbb{E}_t\left[ \int_t^\infty e^{-\rho (s- t)} C_{s} ds \right]
\end{equation*}
subject to
\begin{equation*}
    d W_t = \left[ r W_t + \pi \alpha_t W_t + I_t - C_t  \right] dt + \sigma \alpha_t W_t d Z_t.
\end{equation*}

The HJB equation for this problem is
\begin{equation*}
    \rho V = \max_{C, \alpha} \frac{C^{1-\gamma}}{1-\gamma} + V_W\left[ r W + \pi \alpha W + I  - C\right] + V_I \mu I+ \frac{1}{2} V_{WW} W^2 \alpha^2 \sigma^2 + V_{WI} W  I \alpha \sigma^2 + \frac{1}{2} V_{II} I^2 \sigma^2
\end{equation*}

The optimal consumption and portfolio share satisfy the conditions
\begin{equation*}
 C = V_W^{-\frac{1}{\gamma}}, \qquad     \alpha = -\frac{V_{W} \pi}{V_{WW} W \sigma^2}   - \frac{V_{WI}I}{V_{WW} W}
\end{equation*}

The value function for this problem is given by
\begin{equation*}
    V(W,I) = A \frac{\left( W + H(I) \right)^{1-\gamma}}{1-\gamma},
\end{equation*}
where $H(I) = \frac{I}{r+ \pi - \mu}$ is the investor's human wealth.

The portfolio share is given by
\begin{equation*}
    \frac{\alpha W}{W+H} = \frac{\pi}{\gamma\sigma^2} - \frac{H}{W+H} 
\end{equation*}

\begin{equation*}
    \frac{P Q}{W} = \frac{1+\frac{H}{W}}{\gamma \sigma^2} \pi - \frac{H}{W} 
\end{equation*}


Using $W = \psi (W+H)$ and $H = (1-\psi) (W+H)$ and denoting the quantity demanded by the investor is given by
\begin{equation*}
    \psi \frac{P Q}{W} = \frac{\pi}{\gamma\sigma^2} - (1-\psi)
\end{equation*}

In equilibrium, we have that $P = W$, $Q = 1$, and $\pi = \gamma \sigma^2$. The impact of a small (permanent) change in the price on the quantity demanded is given by
\begin{equation*}
    q = - \frac{1}{\psi \pi} \frac{D_t}{P_t} p
\end{equation*}
where $p = \frac{dP}{P}$. 

\begin{equation*}
    \frac{d q}{q} = (1-\psi^{-1}) \frac{1-\alpha_p}{\omega_A} \frac{\gamma \sigma^2}{c^*}  \frac{F}{W}
\end{equation*}

\newpage
\singlespacing
%\bibliographystyle{jf}
\phantomsection\addcontentsline{toc}{section}{\refname}
\clearpage
%{\small
\bibliography{references.bib}
\end{document}